\documentclass[12pt,a4paper]{article}

\usepackage[T1]{fontenc}
\usepackage[utf8]{inputenc}
\usepackage[catalan]{babel}
\usepackage[margin=2.5cm]{geometry}
\usepackage{lmodern}
\usepackage{amsmath}
\usepackage{booktabs}
\usepackage{array}
\usepackage{tabularx}
\usepackage{xcolor}
\usepackage{float}
\usepackage{fancyhdr}
\usepackage{setspace}
\usepackage{tcolorbox}
\tcbuselibrary{skins}
\usepackage{enumitem}

\setlength{\headheight}{14.5pt}
\pagestyle{fancy}
\fancyhf{}
\fancyhead[L]{Física 4t ESO --- MRU per Fases}
\fancyhead[R]{Cinemàtica}
\fancyfoot[C]{\thepage}
\setlist[itemize]{leftmargin=*, itemsep=0.3em}

\begin{document}

\begin{center}
  {\Large\bfseries Resolució de Problema de Cinemàtica (MRU per Fases)}\\[0.4em]
  {\large Física 4t ESO}\\[0.2em]
  {\normalsize Oriol Farrés --- XC Gradient}
\end{center}
\vspace{0.5em}
\hrule
\vspace{1em}

\begin{tcolorbox}[title=Enunciat,colframe=blue!60!black,colback=blue!5,sharp corners=south]
Un alumne es lleva a les 06:00:00 i surt de Barcelona (BCN) a les 06:32:00. Realitza un viatge en tres fases fins a Madrid (MAD) passant per Guadalajara (GDL). Les dades del viatge són les següents:
\begin{itemize}
  \item Hora de llevar-se: 06:00:00
  \item Hora de sortida (BCN): 06:32:00
  \item Fase 1 (BCN $\to$ GDL): $\Delta x_1 = 350\,\text{km}$ a $v_1 = 40\,\text{m/s}$
  \item Fase 2 (Parada a GDL): $\Delta t_2 = 1\,\text{h}\ 23\,\text{min}$
  \item Fase 3 (GDL $\to$ MAD): $\Delta x_3 = 250\,\text{km}$ a $v_3 = 45\,\text{m/s}$
\end{itemize}
Calculeu el temps total de viatge des que l'alumne es lleva i determineu l'hora d'arribada a Madrid.
\end{tcolorbox}

\section{Nomenclatura Estàndard}

\begin{table}[H]
\centering
\begin{tabularx}{\textwidth}{@{}>{\raggedright\arraybackslash}p{0.18\textwidth}>{\centering\arraybackslash}p{0.14\textwidth}>{\raggedright\arraybackslash}X>{\centering\arraybackslash}p{0.18\textwidth}@{}}
\toprule
Concepte & Símbol & Definició & Unitat (SI) \\
\midrule
instant & $t$ & Moment puntual de temps & s \\
posició & $x$ & Localització sobre l'eix & m \\
tram & $\Delta x$ & Distància recorreguda en una fase & m \\
interval & $\Delta t$ & Durada d'una fase & s \\
\bottomrule
\end{tabularx}
\end{table}

\section{Resolució Pas a Pas}

\subsection{Preparació de dades (Conversió d'unitats)}

\begin{itemize}
  \item $\Delta x_1 = 350\,\text{km} \times 1000 = 350\,000\,\text{m}$
  \item $\Delta x_3 = 250\,\text{km} \times 1000 = 250\,000\,\text{m}$
  \item $\Delta t_2 = 1\,\text{h}\ 23\,\text{min} = 3600\,\text{s} + 23 \times 60\,\text{s} = 3600\,\text{s} + 1380\,\text{s} = 4\,980\,\text{s}$
  \item $\Delta t_0 = 32\,\text{min} = 32 \times 60\,\text{s} = 1\,920\,\text{s}$
\end{itemize}

\subsection{Càlcul d'Intervals de Temps}

Apliquem la definició de velocitat mitjana per a MRU:

\begin{equation}
v = \frac{\Delta x}{\Delta t} \implies \Delta t = \frac{\Delta x}{v}
\label{eq:mrv}
\end{equation}

Apliquem a cada fase:

\begin{align*}
\Delta t_1 &= \frac{350\,000\,\text{m}}{40\,\text{m/s}} = 8\,750\,\text{s} \\[6pt]
\Delta t_3 &= \frac{250\,000\,\text{m}}{45\,\text{m/s}} \approx 5\,555{,}56\,\text{s}
\end{align*}

\subsection{Temps Total Acumulat}

El temps total des que l'alumne es lleva és la suma de tots els intervals:

\begin{align}
T_{\text{total}} &= \Delta t_0 + \Delta t_1 + \Delta t_2 + \Delta t_3 \\
                 &= 1\,920\,\text{s} + 8\,750\,\text{s} + 4\,980\,\text{s} + 5\,555{,}56\,\text{s} \\
                 &= \mathbf{21\,205{,}56\,\text{s}}
\label{eq:ttot}
\end{align}

\subsection{Conversió a Hores, Minuts i Segons}

Convertim el resultat a format horari:

\begin{align*}
21\,205{,}56\,\text{s} \div 3600 &= 5\,\text{h} \quad \text{(resta } 21\,205{,}56 - 18\,000 = 3\,205{,}56\,\text{s)} \\
3\,205{,}56\,\text{s} \div 60    &= 53\,\text{min} \quad \text{(resta } 3\,205{,}56 - 3\,180 = 25{,}56\,\text{s} \approx 25\,\text{s)} \\
\Rightarrow T_{\text{total}}     &= 5\,\text{h}\ 53\,\text{min}\ 25\,\text{s}
\end{align*}

Sumant al moment en què l'alumne es lleva (06:00:00):

\begin{equation}
t_{\text{arribada}} = 06{:}00{:}00 + 5\,\text{h}\ 53\,\text{min}\ 25\,\text{s} = \mathbf{11{:}53{:}25}
\label{eq:arriba}
\end{equation}

\section{Resultats}

\begin{tcolorbox}[title=Resultats Finals,colframe=green!50!black,colback=green!8,sharp corners=south]
\textbf{Temps total (des que es lleva):}
\[
T_{\text{total}} = \mathbf{21\,205{,}56\,\text{s}} = \mathbf{5\,\text{h}\ 53\,\text{min}\ 25\,\text{s}}
\]
\textbf{Hora d'arribada a Madrid:}
\[
t_{\text{arribada}} = \mathbf{11{:}53{:}25}
\]
\end{tcolorbox}

\end{document}
